\documentclass[9pt]{extarticle}
\usepackage[T1]{fontenc}
\usepackage{amsmath,amssymb,amsthm,mathtools,microtype}
\usepackage[dvipsnames]{xcolor}
\usepackage[a4paper,margin=0.68in]{geometry}
\usepackage{titlesec}
\usepackage[colorlinks=true,linkcolor=MidnightBlue,citecolor=MidnightBlue,urlcolor=MidnightBlue]{hyperref}
\titlespacing*{\section}{0pt}{1.15ex plus .3ex minus .2ex}{0.5ex}
\newtheorem{theorem}{Theorem}
\newtheorem{lemma}[theorem]{Lemma}
\newcommand{\R}{\mathbb R}
\newcommand{\Z}{\mathbb Z}
\newcommand{\N}{\mathbb N}
\newcommand{\supp}{\operatorname{supp}}
\newcommand{\selfs}{\mathrm{selfs}}
\newcommand{\unif}{\mathrm{unif}}
\newcommand{\M}{M}

\title{The $p\ne q$ selfsimilar Besov multiplier question:\\
\large a sharp high-smoothness phase split and a counterexample}
\author{Autonomous solve-first investigation --- lane 12}
\date{13 August 2026}

\begin{document}
\maketitle

\begin{abstract}
Schneider--Vyb\'iral ask whether their identity
$\M(B^s_{p,p})=B^s_{p,p,\selfs}$ for $0<p\leq1$ can extend to
$p\ne q$.  On the line, the answer is negative for every
$0<q<p\leq1$ and $s\geq1/p$.  We construct one wavelet at each scale,
placed near $2^j$, with its coefficient normalized in $B^s_{p,q}$.
Every rescaled unit localization sees at most one component, so the
infinite sum belongs to $B^s_{p,q,\selfs}$; separated cutoff tests grow
like $N^{1/p}$ before multiplication and $N^{1/q}$ afterwards, so it is
not a multiplier.  Combined with Nguyen--Sickel's multiplier theorem,
this gives, for $s>1/p$, equality when $p\leq q$ and strict containment
when $q<p$.
\end{abstract}

\section{Question and result}

Fix the cutoff $\theta\in C_c^\infty(\R)$ used in the source definition,
and write
\[
 \|f\|_{B^s_{p,q,\selfs}}
 =\sup_{J\in\N_0,\,\ell\in\Z}
   \|\theta(\cdot-\ell)f(2^{-J}\cdot)\|_{B^s_{p,q}}.
\]
Theorem~5.4(ii) of Schneider--Vyb\'iral \cite{SV} proves
$\M(B^s_{p,p})=B^s_{p,p,\selfs}$ for $0<p\leq1$; the following remark
asks whether this can be generalized to $p\ne q$.

\begin{theorem}[High-smoothness phase split]\label{thm:main}
On $\R$ the following hold.
\begin{enumerate}
 \item If $0<p\leq q\leq\infty$ and $s>1/p$, then
 \[
   \M(B^s_{p,q})=B^s_{p,q,\selfs}=B^s_{p,q,\unif}.
 \]
 \item If $0<q<p\leq1$ and $s\geq1/p$, then
 \[
   \M(B^s_{p,q})\subsetneq B^s_{p,q,\selfs}.
 \]
 In particular, the proposed $p\ne q$ generalization is false, including
 at the endpoint $s=1/p$.
\end{enumerate}
\end{theorem}

Part (1) is short.  Nguyen--Sickel's Theorem~1.2 \cite{NS} says that
$\M(B^s_{p,q})=B^s_{p,q,\unif}$ in this range.  Taking $J=0$ in the
selfsimilar norm (and using the standard equivalence of smooth uniform
localizations) gives
$B^s_{p,q,\selfs}\hookrightarrow B^s_{p,q,\unif}$, while
Theorem~5.4(i) of \cite{SV} gives
$\M(B^s_{p,q})\hookrightarrow B^s_{p,q,\selfs}$.  The three spaces
therefore coincide.  We prove part (2) explicitly.

\section{A dilation-uniform separated wavelet sum}

Choose a compactly supported orthonormal wavelet $w$ with regularity and
vanishing moments beyond $s$, and put
\[
 w_{j,k}(x)=2^{j/2}w(2^jx-k),\qquad
 k_j=4^j,\qquad c_j=2^{-j}k_j=2^j.
\]
Let $\supp w\subset[-T,T]$, set
\begin{equation}\label{eq:m}
 A=s+\frac12-\frac1p,\qquad a_j=2^{-jA},
 \qquad m=\sum_{j\geq j_0}a_jw_{j,k_j}.
\end{equation}
and take $j_0$ sufficiently large below.  The supports
\[
 I_j=\supp w_{j,k_j}\subset
 [c_j-T2^{-j},c_j+T2^{-j}]
\]
escape to infinity and are disjoint, so \eqref{eq:m} is a locally finite bounded
function; indeed $a_j2^{j/2}=2^{-j(s-1/p)}\leq1$.

\begin{lemma}\label{lem:selfs}
If $s\geq1/p$, then $m\in B^s_{p,q,\selfs}(\R)$.
\end{lemma}

\begin{proof}
After the dilation in the selfsimilar norm, the $j$th support is
\begin{equation}\label{eq:support}
 2^JI_j\subset
 [2^{J+j}-T2^{J-j},\,2^{J+j}+T2^{J-j}].
\end{equation}
The gap between consecutive intervals before dilation is
$2^j-T(2^{-j}+2^{-j-1})$.  Choose $j_0$ so that this exceeds twice the
diameter of $\supp\theta$.  Since dilation only increases the gaps, every
$\theta(\cdot-\ell)$ meets at most one interval in \eqref{eq:support}, uniformly in
$J$ and $\ell$.

Suppose first that the surviving index satisfies $j\geq J$.  Since
\[
 w_{j,k_j}(2^{-J}x)=2^{J/2}w_{j-J,k_j}(x),
\]
the wavelet characterization of $B^s_{p,q}$ and uniform boundedness of
multiplication by translates of $\theta$ give
\begin{align}
 \|\theta(\cdot-\ell)a_jw_{j,k_j}(2^{-J}\cdot)\|_{B^s_{p,q}}
 &\lesssim a_j2^{J/2}2^{(j-J)(s+1/2-1/p)}\\
 &=2^{-J(s-1/p)}\leq1. \label{eq:fine}
\end{align}

If $j<J$, the component has frequency below one.  For every integer
$0\leq r\leq K$, where $K>s$, Leibniz' rule yields
\begin{equation}\label{eq:broad}
 \big\|D^r[\theta(\cdot-\ell)a_j2^{j/2}
 w(2^{j-J}\cdot-k_j)]\big\|_\infty
 \lesssim a_j2^{j/2}=2^{-j(s-1/p)}\leq1.
\end{equation}
Here every derivative falling on $w$ contributes
$2^{j-J}\leq1$, and derivatives falling on the fixed cutoff are harmless.
The function in \eqref{eq:broad} has uniformly bounded support, so the standard
$C_c^K\hookrightarrow B^s_{p,q}$ estimate gives a uniform Besov bound.
Equations \eqref{eq:fine}--\eqref{eq:broad}, together with the
one-component property, prove the
claim.
\end{proof}

\section{The multiplier obstruction}

Choose $\eta\in C_c^\infty(\R)$ equal to one on $[-T,T]$.  Enlarging
$j_0$ if necessary, all translates $\eta(\cdot-c_j)$ and all of their
supports enlarged by the translations occurring in an order-$r>s$
finite difference are mutually disjoint.  For $N\geq1$, define
\[
 g_N=\sum_{j=j_0}^{j_0+N-1}\eta(\cdot-c_j).
\]
The difference characterization of Besov spaces then gives, with constants
independent of $N$,
\begin{equation}\label{eq:input}
 \|g_N\|_{B^s_{p,q}}\lesssim N^{1/p}.
\end{equation}
Indeed the $p$th powers of the $L_p$ norms add for the separated
translates, both for the functions and for every finite difference.

The cutoff was chosen so that $g_N$ selects exactly the indicated wavelet
components of $m$.  Hence
\[
 mg_N=\sum_{j=j_0}^{j_0+N-1}a_jw_{j,k_j}.
\]
The wavelet characterization and the normalization of $a_j$ imply
\begin{equation}\label{eq:output}
 \|mg_N\|_{B^s_{p,q}}\asymp
 \left(\sum_{j=j_0}^{j_0+N-1}
 [2^{j(s+1/2-1/p)}a_j]^q\right)^{1/q}\asymp N^{1/q}.
\end{equation}
If $m$ were a pointwise multiplier,
\eqref{eq:input}--\eqref{eq:output} would force
$N^{1/q}\lesssim N^{1/p}$ for all $N$, impossible because $q<p$.
Together with Lemma~\ref{lem:selfs}, this proves part (2) of
Theorem~\ref{thm:main}.

\section{Scope and verification notes}

The construction is one-dimensional, which already disproves a general
extension on $\R^n$.  It works at $s=1/p$: the right sides of
\eqref{eq:fine}--\eqref{eq:broad}
then remain constant rather than decay.  The hypothesis $s\geq1/p$ also
places the source's difference-defined Besov space safely in the range
where the standard Fourier/wavelet description agrees with it.

Nguyen--Sickel use the same separated translation indices in proving the
necessity direction for localized multiplier classes.  The additional
step here is the uniform audit under every dilation in the 2012
selfsimilar norm, especially the broad-bump case $j<J$.  Exact-phrase,
arXiv-id, and selfsimilar-multiplier searches located no explicit answer
to the source question.  This is a bounded novelty check, not a priority
claim.

\paragraph{Human-review recommendation.}
Accept as a candidate full negative answer in the stated parameter range.
The key checks are the one-component property after arbitrary source
dilations, the endpoint estimate \eqref{eq:broad}, and the separated
finite-difference bound \eqref{eq:input}.

\begin{thebibliography}{9}
\bibitem{SV}
C.~Schneider and J.~Vyb\'iral,
\emph{Non-smooth atomic decompositions, traces on Lipschitz domains, and
pointwise multipliers in function spaces},
\href{https://arxiv.org/abs/1201.2280}{arXiv:1201.2280} (2012),
Theorem~5.4 and the following remark.

\bibitem{NS}
V.~K.~Nguyen and W.~Sickel,
\emph{On a Problem of Jaak Peetre Concerning Pointwise Multipliers of
Besov Spaces},
\href{https://arxiv.org/abs/1703.03246}{arXiv:1703.03246} (2017),
Theorems~1.2, 1.6 and Proposition~3.10.
\end{thebibliography}

\end{document}
